\documentclass{article}
\usepackage{graphicx} % Required for inserting images
\usepackage{dsfont}
\usepackage{amssymb}  % For set symbols like \mathbb{N}
\usepackage{amsmath}  % For advanced math features
\usepackage{enumitem} % For customizing lists 
\usepackage{amsfonts}
\usepackage{amsthm}
\usepackage{tikz}
\usepackage{mathtools}
\setlength{\parindent}{0pt}
\begin{document}
\title{ASP Parallel Planning}
\author{Levi Klein}
\date{13.11.2025}
\maketitle
\section{Definitions}
A STRIPS-like (multivalued) \textit{planning task} according to blablabla is a 4-tuple $\langle F, s_0, s_*, \mathcal O \rangle$  with
\begin{itemize}[leftmargin=*]
    \item $F$ is a finite set of \emph{fluents} (variables that may change between states). Each $x\in F$ has an associated finite domain $x^d$ of possible values.
    \item $s_0$ is the initial state, i.e. a total function that maps every $x\in F$ to a value in its domain $x^d$.
    \item $s_*$ is a partial state specifying the goal conditions.
    \item $\mathcal O$ is a finite set of operators (actions) over $F$. Each action $a$ has the form $a=\langle a^c,a^e\rangle$, where $a^c$ and $a^e$ are partial states denoting the precondition and postcondition of $a$.
\end{itemize}
\end{document}